\documentclass[11pt]{article}
\usepackage[a4paper,margin=22mm]{geometry}
\usepackage{fontspec}
\setmainfont{DejaVu Serif}
\setsansfont{DejaVu Sans}
\usepackage{microtype}
\usepackage{xcolor}
\usepackage{hyperref}
\usepackage{enumitem}
\usepackage{parskip}
\definecolor{Accent}{HTML}{971D32}
\definecolor{Muted}{HTML}{666666}
\hypersetup{colorlinks=true,urlcolor=Accent,linkcolor=Accent}
\setlist{leftmargin=1.6em,itemsep=0.3em,topsep=0.35em}
\setcounter{tocdepth}{2}
\renewcommand{\contentsname}{Contents}
\newcommand{\sectionrule}{\par\vspace{-0.5em}\noindent\textcolor{Accent}{\rule{\linewidth}{0.6pt}}\par}
\begin{document}

\begin{center}
{\sffamily\bfseries\color{Accent} TOEFL WORD OF THE DAY\par}
\vspace{8mm}
{\fontsize{42}{48}\selectfont\bfseries denote\par}
\vspace{2mm}
{\Large\sffamily\color{Accent} /dɪˈnoʊt/\par}
\vspace{2mm}
{\small\sffamily Local audio phrase: denote\par}
{\small\sffamily Audio file: audios/denote.mp3\par}
\vspace{2mm}
{\large\sffamily\color{Muted} transitive verb\par}
\vspace{5mm}
{\sffamily mean or represent \textbullet\ show or indicate\par}
\vspace{6mm}
{\large To denote something is to represent a meaning or to show that a particular condition exists.\par}
\vspace{6mm}
\begin{minipage}{0.84\textwidth}
\itshape “Different colors on the map denote different types of land use.”
\end{minipage}\par
\vspace{10mm}
\textbf{Date:} August 3rd 2026\quad\textbullet\quad \textbf{Level:} B1--mid-B2 TOEFL
\vfill
Created and compiled by \textbf{Amir Kooshky}\par
\vspace{2mm}
\href{https://t.me/KooshkyTOEFL}{Telegram Channel}\quad\textbullet\quad
\href{https://www.instagram.com/kooshkytoefl}{Instagram}\quad\textbullet\quad
\href{https://t.me/KooshkyTOEFL\_pv}{Personal Telegram}
\end{center}
\newpage
\tableofcontents
\newpage

\section{Meanings and usage}
\sectionrule
\subsection{Meaning 1: Mean or represent}
\textbf{Part of speech:} transitive verb

\textbf{IPA:} /dɪˈnoʊt/

\textbf{Pronunciation phrase:} denote

\textbf{Local audio file:} audios/denote.mp3

\textbf{Register:} formal; common in academic, mathematical, scientific, linguistic, and technical writing

\textbf{Definition:} to have a particular meaning or to represent a particular idea, object, or category

\textbf{In simpler words:} mean or stand for

\textbf{Typical pattern:} word, symbol, sign, or label + denotes + meaning or category

\subsubsection{Natural examples}
\begin{enumerate}
  \item In mathematics, the symbol π denotes the ratio of a circle's circumference to its diameter.
  \item The term renewable energy denotes energy obtained from sources that can naturally be replaced.
  \item Different colors on the map denote different types of land use.
  \item In the diagram, dotted lines denote proposed transportation routes.
  \item The letter N is used to denote the total number of participants.
  \item In this article, the word community denotes a group of people who share a common location or identity.
  \item Each code denotes a different category of archaeological evidence.
\end{enumerate}

\subsubsection{Nuance and implication}
\textbf{Central nuance:} The word describes a direct or established connection between a sign, term, symbol, or label and what it represents.

\textbf{Usual implication:} The meaning is often defined by a system, convention, author, or field of study rather than guessed from context.

\textbf{Compared with connote:} Denote refers to the direct or basic meaning of a word or sign. Connote refers to the additional ideas, feelings, or associations that it suggests.

\subsubsection{Grammar notes}
\textbf{How to use it:} This verb normally needs the meaning, object, or category being represented after it.

\textbf{What can follow it:} It is often followed by a noun or noun phrase that explains what a word, symbol, color, number, or sign represents.

\textbf{Position in a sentence:} The word, symbol, sign, or label usually comes before denote, and its meaning usually comes after it.

\subsubsection{Useful patterns}
\textbf{Pattern:} symbol, letter, or number + denotes + meaning

\textbf{Use:} Use this when explaining what a symbol represents.

\textbf{Example:} The letter V denotes velocity in this equation.

\textbf{Pattern:} term or word + denotes + concept

\textbf{Use:} Use this when defining how a word is used in a particular text or field.

\textbf{Example:} In economics, the term capital denotes resources used to produce goods and services.

\textbf{Pattern:} be denoted by + symbol or label

\textbf{Use:} Use the passive form when the represented idea is the main focus.

\textbf{Example:} Temperature is denoted by the letter T in the formula.

\textbf{Pattern:} use + symbol + to denote + noun

\textbf{Use:} Use this when explaining the purpose of a symbol or label.

\textbf{Example:} The researchers used circles to denote urban locations.

\subsubsection{Common wrong usage}
\paragraph{Correction 1}
\textbf{Wrong:} The symbol denotes to the total population.

\textbf{Correct:} The symbol denotes the total population.

\textbf{Why:} Denote takes its object directly and is not followed by to.

\paragraph{Correction 2}
\textbf{Wrong:} The letter is denoted the average value.

\textbf{Correct:} The average value is denoted by the letter.

\textbf{Why:} In the passive form, use be denoted by to identify the symbol.

\paragraph{Correction 3}
\textbf{Wrong:} The word denotes about social change.

\textbf{Correct:} The word denotes social change.

\textbf{Why:} Do not use about after denote.

\subsection{Meaning 2: Show or indicate}
\textbf{Part of speech:} transitive verb

\textbf{IPA:} /dɪˈnoʊt/

\textbf{Pronunciation phrase:} denote

\textbf{Local audio file:} audios/denote.mp3

\textbf{Register:} formal; common in academic, scientific, medical, and analytical writing

\textbf{Definition:} to show or indicate that a particular condition, quality, or situation exists

\textbf{In simpler words:} show that something is present or true

\textbf{Typical pattern:} feature, result, behavior, or measurement + denotes + condition or quality

\subsubsection{Natural examples}
\begin{enumerate}
  \item A sudden drop in temperature may denote an approaching storm.
  \item In this grading system, a score above 90 denotes excellent performance.
  \item The presence of these chemicals may denote contamination of the water supply.
  \item Dark areas in the image denote regions of lower temperature.
  \item Frequent interruptions may denote a lack of preparation.
  \item An increase in plant diversity often denotes a healthier ecosystem.
  \item The change in color denotes that the chemical reaction is complete.
\end{enumerate}

\subsubsection{Nuance and implication}
\textbf{Central nuance:} The word presents one observable fact as a sign of another condition or quality.

\textbf{Usual implication:} The connection is often based on an accepted rule, classification system, or interpretation of evidence.

\textbf{Compared with cause:} If one thing denotes another, it shows or represents it. This does not necessarily mean that it causes it.

\subsubsection{Grammar notes}
\textbf{How to use it:} This verb normally needs the condition, quality, or situation being indicated after it.

\textbf{What can follow it:} It is commonly followed by nouns such as change, difference, category, status, quality, presence, absence, success, or danger.

\textbf{Position in a sentence:} The observable sign, measurement, behavior, or feature normally comes before the verb.

\subsubsection{Useful patterns}
\textbf{Pattern:} feature or result + denotes + condition

\textbf{Use:} Use this when an observable feature indicates a particular condition.

\textbf{Example:} A high concentration of salt may denote severe water loss.

\textbf{Pattern:} color or symbol + denotes + category or status

\textbf{Use:} Use this when a visual feature identifies a category or condition.

\textbf{Example:} A red label denotes a high level of risk.

\textbf{Pattern:} be used to denote + condition or category

\textbf{Use:} Use this when a sign has been assigned a particular meaning.

\textbf{Example:} A flashing light is used to denote an emergency.

\textbf{Pattern:} may denote + condition

\textbf{Use:} Use may when the sign suggests a condition but does not prove it with certainty.

\textbf{Example:} Persistent fatigue may denote an underlying health problem.

\subsubsection{Common wrong usage}
\paragraph{Correction 1}
\textbf{Wrong:} The red light denotes that of danger.

\textbf{Correct:} The red light denotes danger.

\textbf{Why:} Denote can be followed directly by the condition being indicated.

\paragraph{Correction 2}
\textbf{Wrong:} The result denotes us that the theory is incorrect.

\textbf{Correct:} The result denotes that the theory is incorrect.

\textbf{Why:} Do not place a person between denote and a clause beginning with that.

\paragraph{Correction 3}
\textbf{Wrong:} The symptom denotes from a serious illness.

\textbf{Correct:} The symptom may denote a serious illness.

\textbf{Why:} Denote takes its object directly and is not followed by from.

\section{Word family}
\sectionrule
\subsection{denotation}
\textbf{Part of speech:} countable or uncountable noun

\textbf{IPA:} /ˌdiːnoʊˈteɪʃən/

\textbf{Definition:} the direct or basic meaning of a word, sign, or symbol

\textbf{Relationship to the headword:} This noun refers to the direct meaning that a word or sign denotes.

\textbf{Grammar pattern:} the denotation of + word, sign, or symbol

\textbf{Usage note:} It is especially common in linguistics, literary analysis, and discussions of meaning.

\subsubsection{Examples}
\begin{enumerate}
  \item The denotation of the word home is the place where someone lives.
  \item Words with similar denotations may have very different emotional associations.
\end{enumerate}

\section{Common collocations}
\sectionrule
\subsection{denote a difference}
\textbf{Applies to:} Show or indicate

\textbf{Meaning:} show that two things belong to different groups or have different qualities

\textbf{Example:} The two line styles denote a difference in ownership.

\subsection{denote a change}
\textbf{Applies to:} Show or indicate

\textbf{Meaning:} indicate that a condition has changed

\textbf{Example:} A shift in color may denote a change in acidity.

\subsection{denote a category}
\textbf{Applies to:} Mean or represent

\textbf{Meaning:} represent one particular group in a classification system

\textbf{Example:} Each number denotes a separate category of response.

\subsection{denote the presence of}
\textbf{Applies to:} Show or indicate

\textbf{Meaning:} show that something exists in a place or sample

\textbf{Pattern:} denote the presence of + substance or condition

\textbf{Example:} Cloudy water may denote the presence of suspended particles.

\subsection{denote the absence of}
\textbf{Applies to:} Show or indicate

\textbf{Meaning:} show that something is not present

\textbf{Pattern:} denote the absence of + noun

\textbf{Example:} A blank field denotes the absence of available data.

\subsection{be denoted by}
\textbf{Applies to:} all senses

\textbf{Meaning:} be represented or indicated by a particular symbol, sign, or feature

\textbf{Pattern:} noun + be denoted by + symbol, sign, or label

\textbf{Example:} The control group is denoted by a square in the graph.

\subsection{use something to denote}
\textbf{Applies to:} all senses

\textbf{Meaning:} give a symbol, color, or sign a particular meaning

\textbf{Pattern:} use + symbol or feature + to denote + meaning

\textbf{Example:} The map uses blue shading to denote areas at risk of flooding.

\subsection{the term denotes}
\textbf{Applies to:} Mean or represent

\textbf{Meaning:} the term has a particular defined meaning

\textbf{Pattern:} the term + denotes + concept

\textbf{Example:} In biology, the term adaptation denotes a characteristic that improves survival or reproduction.

\subsection{the symbol denotes}
\textbf{Applies to:} Mean or represent

\textbf{Meaning:} the symbol represents a particular quantity or idea

\textbf{Pattern:} the symbol + denotes + quantity or concept

\textbf{Example:} In the equation, the symbol R denotes resistance.

\section{Synonyms and antonyms}
\sectionrule
\subsection{Close synonyms}
\subsubsection{represent}
\textbf{Part of speech:} transitive verb

\textbf{Applies to:} Mean or represent

\textbf{Shared meaning:} Both words can describe one thing standing for another.

\textbf{Difference:} Represent is broader and can describe pictures, people, organizations, or symbols standing for something. Denote is more formal and focuses on a specific established meaning.

\textbf{Example:} The dove represents peace in many cultures.

\subsubsection{signify}
\textbf{Part of speech:} transitive verb

\textbf{Applies to:} all senses

\textbf{Shared meaning:} Both words can mean represent, mean, or indicate.

\textbf{Difference:} Signify is very close to denote but may place slightly more emphasis on importance or interpretation. Denote often sounds more precise and technical.

\textbf{Example:} The bell signifies the end of the examination.

\subsubsection{indicate}
\textbf{Part of speech:} transitive verb

\textbf{Applies to:} Show or indicate

\textbf{Shared meaning:} Both words can mean show that a condition or fact exists.

\textbf{Difference:} Indicate is more common and can mean point toward, suggest, or provide evidence of something. Denote often suggests a clearer or more established connection.

\textbf{Example:} The results indicate that the treatment was effective.

\subsubsection{refer to}
\textbf{Part of speech:} verb phrase

\textbf{Applies to:} Mean or represent

\textbf{Shared meaning:} Both expressions can connect a term with the thing or concept it identifies.

\textbf{Difference:} Refer to identifies the person, thing, or idea being discussed. Denote emphasizes the meaning assigned to a word, symbol, or sign.

\textbf{Example:} In this chapter, the term youth refers to people between the ages of 15 and 24.

\section{Words learners may confuse}
\sectionrule
\subsection{connote}
\textbf{Part of speech:} transitive verb

\textbf{Why learners confuse them:} Both words describe meaning and are frequently contrasted in language and literature courses.

\textbf{Key difference:} Denote means have a direct or basic meaning. Connote means suggest additional feelings, ideas, or associations beyond that basic meaning.

\textbf{denote:} The word snake denotes a long, legless reptile.

\textbf{connote:} The word snake may connote danger or dishonesty.

\subsection{designate}
\textbf{Part of speech:} transitive verb

\textbf{Why learners confuse them:} The words are formal, have similar forms, and can both involve labels or categories.

\textbf{Key difference:} Denote means represent or indicate something. Designate usually means officially choose, name, or assign something for a particular purpose.

\textbf{denote:} A green circle denotes an open station.

\textbf{designate:} The city designated the building as a protected historical site.

\section{Etymology}
\sectionrule
\textbf{Origin:} Denote comes from Latin denotare, meaning to mark out or point out.

\section{Practice exercises}
\sectionrule
\subsection{Word family choice}
Choose the best member of the word family for each blank.

\begin{enumerate}
  \item The literal \_\_\_\_\_ of the word rose is a type of flowering plant.
  \item In this formula, the letter M is used to \_\_\_\_\_ mass.
\end{enumerate}

\subsection{Correct or incorrect?}
Decide whether each sentence is correct. Correct the inaccurate ones.

\begin{enumerate}
  \item The symbol denotes the total amount of energy produced.
  \item The red line denotes to the national border.
  \item A high reading may denote the presence of pollution.
  \item The word inexpensive denotes feelings of poor quality.
\end{enumerate}

\subsection{Collocation practice}
Complete each sentence with a natural collocation.

\begin{enumerate}
  \item On this map, a red triangle is used to \_\_\_\_\_ an active volcano.
  \item The experimental group is denoted \_\_\_\_\_ a solid black line.
  \item A sudden increase in pressure may denote the \_\_\_\_\_ of a blockage.
  \item In the table, different letters denote separate \_\_\_\_\_.
\end{enumerate}

\subsection{Complete answer key}
\subsubsection{Word family choice}
\begin{enumerate}
  \item \textbf{denotation.} The literal denotation of the word rose is a type of flowering plant. \emph{Use the noun denotation for the direct or basic meaning of a word.}
  \item \textbf{denote.} In this formula, the letter M is used to denote mass. \emph{Use the base verb denote after to.}
\end{enumerate}

\subsubsection{Correct or incorrect?}
\begin{enumerate}
  \item \textbf{Correct} \emph{Denote correctly takes the represented quantity directly as its object.}
  \item \textbf{Incorrect}. The red line denotes the national border. \emph{Do not use to after denote.}
  \item \textbf{Correct} \emph{Denote the presence of is a natural formal expression.}
  \item \textbf{Incorrect}. The word inexpensive may connote poor quality. \emph{Use connote for an additional association rather than a direct dictionary meaning.}
\end{enumerate}

\subsubsection{Collocation practice}
\begin{enumerate}
  \item \textbf{denote.} On this map, a red triangle is used to denote an active volcano. \emph{Use something to denote means give a sign or symbol a particular meaning.}
  \item \textbf{by.} The experimental group is denoted by a solid black line. \emph{Use be denoted by to identify the sign or symbol representing something.}
  \item \textbf{presence.} A sudden increase in pressure may denote the presence of a blockage. \emph{Denote the presence of means show that something exists.}
  \item \textbf{categories.} In the table, different letters denote separate categories. \emph{Denote a category means represent a particular group in a classification system.}
\end{enumerate}

\vfill
\begin{center}
\small\color{Muted} Created and compiled by \textbf{Amir Kooshky}\\
\href{https://t.me/KooshkyTOEFL}{Telegram Channel}\quad\textbullet\quad
\href{https://www.instagram.com/kooshkytoefl}{Instagram}\quad\textbullet\quad
\href{https://t.me/KooshkyTOEFL\_pv}{Personal Telegram}
\end{center}
\end{document}
